\begin{afterword}
\summaryhead

The author introduces a thought experiment made to break the author's own optimism. The previous post had traced the author's folly to a felt conviction that things would turn out all right in the end.

A God taken seriously, as a child takes it, need not grant every request, but would intervene past some threshold of horror; a God who never intervenes is a belief protected from falsification. In that world no soul need fear annihilation, and even a torture simulated in Conway's Game of Life would be stopped. But the mathematical answer to what given rules and a given starting state would produce is beyond even God's reach. In that what-if world, life may evolve and suffer, a Genghis Khan may torture innocents and die happy, and the fate of millions may turn on the timing of one conception.

The author holds that we live in the what-if world. Many atheists still think that some things are ``not allowed'': that World War II would have happened without Hitler, that liberal democracies will not legalize torture, that technology cannot do more harm than good. Readers who value happiness above all need not dwell on this. Those with something to protect must, since the only fairer world is one we build, starting from this one.

\responsehead

The post's thesis is sound and carefully limited. No physical law guarantees that catastrophe will not happen, and the post presents seeing this as a step before arguing probabilities, not a substitute for it. The Game of Life makes the point concrete. The history it uses is accurate: Hitler's lost chance at architecture, and the general view that Hitler wanted and launched the war. Its comparison of our world with an indifferent one is the comparison Paul Draper made formally in 1989.

The post also explains other people's views. It says the ``clever arguments'' for why the positive-sum trend cannot reverse, or why liberal democracies will not legalize torture, have as their source ``a much deeper belief that such things are not allowed''. This is stated without a hedge; no one is named and no argument is quoted, and the only case given is the author's own former belief. On World War II the post is hedged: it knows of no empirical reason to doubt that events turned on Hitler ``that I happen to have heard of.'' There is a documented dissent, A.~J.~P. Taylor's, argued from diplomatic documents and rejected by its critics on the evidence.

On God, the post takes a side in the problem of evil, and the reader should know the other side. The post holds that a loving God, like a loving parent, would draw the line at some horror, so a God who never intervenes is a belief protected from falsification. The theist reply reported in the Stanford Encyclopedia of Philosophy is that the goods which justify God in permitting an evil may lie beyond our ken. On that reply, God's not intervening in the horrors we know of is neither a God who never intervenes nor a belief built to avoid falsification. William Rowe's answer, that the reply would undercut inductive reasoning in general, is on the post's side.

In short: a careful thought experiment for the sound claim that nothing in physics forbids catastrophe, with an unhedged diagnosis of other people's arguments to the contrary, whose only case is the author's own former belief.
\end{afterword}
